\section{Encoding Dichotomy and ETH Lower Bounds}\label{sec:dichotomy}

The hardness results of Section~\ref{sec:hardness} apply to worst-case instances under the succinct encoding. This section states a dichotomy that separates an explicit-state upper bound from a succinct-encoding lower bound. The dichotomy and ETH reduction chain are machine-verified in Lean 4.

\paragraph{Model note.} Part~1 is an explicit-state upper bound (time polynomial in $|S|$). Part~2 is a succinct-encoding lower bound under ETH (time exponential in $n$). The encodings are defined in Section~\ref{sec:encoding}.

\begin{theorem}[Complexity Dichotomy]\label{thm:dichotomy}
Let $\mathcal{D} = (A, X_1, \ldots, X_n, U)$ be a decision problem with $|S| = N$ states. Let $k^*$ be the size of the minimal sufficient set.

\begin{enumerate}
\item \textbf{Logarithmic case (explicit-state upper bound):} If $k^* = O(\log N)$, then SUFFICIENCY-CHECK is solvable in polynomial time in $N$ under the explicit-state encoding.

\item \textbf{Linear case (succinct lower bound under ETH):} If $k^* = \Omega(n)$, then SUFFICIENCY-CHECK requires time $\Omega(2^{n/c})$ for some constant $c > 0$ under the succinct encoding (assuming ETH).
\end{enumerate}
\end{theorem}

\begin{proof}
\textbf{Part 1 (Logarithmic case):} If $k^* = O(\log N)$, then the number of distinct projections $|S_{I^*}|$ is at most $2^{k^*} = O(N^c)$ for some constant $c$. Under the explicit-state encoding, we enumerate all projections and verify sufficiency in polynomial time.

\textbf{Part 2 (Linear case):} We establish this via an explicit reduction chain from the Exponential Time Hypothesis under the succinct encoding.
\end{proof}

\subsection{The ETH Reduction Chain}\label{sec:eth-chain}

The lower bound in Part 2 of Theorem~\ref{thm:dichotomy} follows from a chain of reductions originating in the Exponential Time Hypothesis. We make this chain explicit.

\begin{definition}[Exponential Time Hypothesis (ETH)]
There exists a constant $\delta > 0$ such that 3-SAT on $n$ variables cannot be solved in time $O(2^{\delta n})$~\cite{impagliazzo2001complexity}.
\end{definition}

The chain proceeds as follows:

\begin{enumerate}
\item \textbf{ETH $\Rightarrow$ 3-SAT requires $2^{\Omega(n)}$:} This is the definition of ETH.

\item \textbf{3-SAT $\leq_p$ TAUTOLOGY:} Given 3-SAT formula $\varphi(x_1, \ldots, x_n)$, define $\psi = \neg\varphi$. Then $\varphi$ is satisfiable iff $\psi$ is not a tautology. This is a linear-time reduction preserving the number of variables.

\item \textbf{TAUTOLOGY requires $2^{\Omega(n)}$ (under ETH):} By the contrapositive of step 2, if TAUTOLOGY is solvable in $o(2^{\delta n})$ time, then 3-SAT is solvable in $o(2^{\delta n})$ time, contradicting ETH.

\item \textbf{TAUTOLOGY $\leq_p$ SUFFICIENCY-CHECK:} This is Theorem~\ref{thm:sufficiency-conp}. Given formula $\varphi(x_1, \ldots, x_n)$, we construct a decision problem where:
\begin{itemize}
\item The empty set $I = \emptyset$ is sufficient iff $\varphi$ is a tautology
\item When $\varphi$ is not a tautology, all $n$ coordinates are relevant
\end{itemize}
The reduction is polynomial-time and preserves the number of coordinates.

\item \textbf{SUFFICIENCY-CHECK requires $2^{\Omega(n)}$ (under ETH):} Combining steps 3 and 4: if SUFFICIENCY-CHECK is solvable in $o(2^{\delta n/c})$ time for some constant $c$, then TAUTOLOGY (and hence 3-SAT) is solvable in subexponential time, contradicting ETH.
\end{enumerate}

\begin{proposition}[Tight Constant]\label{prop:eth-constant}
The reduction in Theorem~\ref{thm:sufficiency-conp} preserves the number of variables up to an additive constant: an $n$-variable formula yields an $(n+1)$-coordinate decision problem. Therefore, the constant $c$ in the $2^{n/c}$ lower bound is asymptotically 1:
\[
\text{SUFFICIENCY-CHECK requires time } \Omega(2^{\delta (n-1)}) = 2^{\Omega(n)} \text{ under ETH}
\]
where $\delta$ is the ETH constant for 3-SAT.
\end{proposition}

\begin{proof}
The TAUTOLOGY reduction (Theorem~\ref{thm:sufficiency-conp}) constructs:
\begin{itemize}
\item State space $S = \{\text{ref}\} \cup \{0,1\}^n$ with $n+1$ coordinates (one extra for the reference state)
\item Query set $I = \emptyset$
\end{itemize}
When $\varphi$ has $n$ variables, the constructed problem has $n+1$ coordinates. The asymptotic lower bound is $2^{\Omega(n)}$ with the same constant $\delta$ from ETH.
\end{proof}

\subsection{Phase Transition}

\begin{corollary}[Phase Transition]\label{cor:phase-transition}
There is a sharp transition between tractable and intractable regimes at the logarithmic scale (with respect to the encodings in Section~\ref{sec:encoding}):
\begin{itemize}
\item If $k^* = O(\log N)$, SUFFICIENCY-CHECK is polynomial in $N$ under the explicit-state encoding
\item If $k^* = \Omega(n)$, SUFFICIENCY-CHECK is exponential in $n$ under ETH in the succinct encoding
\end{itemize}
For Boolean coordinate spaces ($N = 2^n$), the transition is between $k^* = O(\log n)$ (explicit-state tractable) and $k^* = \Omega(n)$ (succinct-encoding intractable).
\end{corollary}

\begin{proof}
The logarithmic case (Part 1 of Theorem~\ref{thm:dichotomy}) gives polynomial time when $k^* = O(\log N)$. More precisely, when $k^* \leq c \log N$ for constant $c$, the algorithm runs in time $O(N^c \cdot \text{poly}(n))$.

The linear case (Part 2) gives exponential time when $k^* = \Omega(n)$.

The transition point is where $2^{k^*} = N^{k^*/\log N}$ stops being polynomial in $N$, i.e., when $k^* = \omega(\log N)$. For Boolean coordinate spaces ($N = 2^n$), this corresponds to the gap between $k^* = O(\log n)$ and $k^* = \Omega(n)$.
\end{proof}

\begin{remark}[Sharpness of Dichotomy]
Under ETH, the lower bound is asymptotically tight in the succinct encoding. The explicit-state upper bound is tight in the sense that it matches enumeration complexity in $N$.
\end{remark}

This dichotomy explains why some domains admit tractable model selection (few relevant variables) while others require heuristics (many relevant variables). In the succinct regime, the reduction chain gives a class-level coNP statement together with a worst-case $2^{\Omega(n)}$ lower bound under the declared ETH assumption.

\paragraph{Bridge theorems.} The ETH chain composes with structural rank characterizations to yield the complete complexity picture. Specifically: \LH{BR1} combines ETH hardness with structural rank to show that problems with maximal srank (all coordinates relevant) inherit the exponential lower bound; \LH{CC1} composes the dichotomy theorem with ETH to give the full classification; \LH{DQ5} combines the TAUTOLOGY reduction with ETH to establish coNP-completeness plus exponential hardness. These bridge theorems connect the complexity cluster (ETH, reductions) to the geometric cluster (srank, sufficiency) in the proof graph.

\begin{remark}[Circuit Model Formalization]
The ETH lower bound is stated in the succinct circuit encoding of Section~\ref{sec:encoding}, where the utility function $U: A \times S \to \mathbb{R}$ is represented by a Boolean circuit computing $\mathbf{1}[U(a,s) > \theta]$ for threshold comparisons. In this model:
\begin{itemize}
\item The input size is the circuit size $m$, not the state space size $|S| = 2^n$
\item A 3-SAT formula with $n$ variables and $c$ clauses yields a circuit of size $O(n + c)$
\item The reduction preserves instance size up to constant factors: $m_{\text{out}} \leq 3 \cdot m_{\text{in}}$
\end{itemize}
This linear size preservation is essential for ETH transfer. In the explicit enumeration model (where $S$ is given as a list), the reduction would blow up the instance size exponentially, precluding ETH-based lower bounds. The circuit model is standard in fine-grained complexity and matches practical representations of decision problems.
\end{remark}
